\begin{table}
\caption{Magnetic and structure-aware physics for the Curie temperature}
\label{tab:table3_magnetic_physics}
\tablefont
\begin{tabularx}{\textwidth}{@{}>{\raggedright\arraybackslash}Xrllll@{}}
\toprule
\textbf{Algorithm} & \textbf{MAE (K)} & \textbf{vs Tier 1} & \textbf{Verdict vs Tier 1} & \textbf{vs Tier 3} & \textbf{Verdict vs Tier 3} \\
\midrule
\multicolumn{6}{@{}l@{}}{\textit{Tier 1 + magnetic-sublattice geometry (10)}} \\
\quad Gradient boosting & 202.6 & +0.3 [$-$13.9, +14.5] & none & -- & -- \\
\quad XGBoost & 203.6 & +5.8 [$-$15.8, +27.4] & none & -- & -- \\
\addlinespace
\multicolumn{6}{@{}l@{}}{\textit{Tier 3 + magnetic-sublattice geometry}} \\
\quad Gradient boosting & 220.2 & -- & -- & $-$3.1 [$-$27.6, +21.4] & none \\
\quad XGBoost & 211.0 & -- & -- & +8.0 [$-$15.8, +31.9] & none \\
\addlinespace
\multicolumn{6}{@{}l@{}}{\textit{Tier 3 + DFT magnetic outputs (4)}} \\
\quad Gradient boosting & 197.6 & +5.3 [$-$29.9, +40.4] & none & +19.6 [$-$3.1, +42.3] & borderline gain \\
\quad XGBoost & 210.7 & $-$1.4 [$-$35.8, +33.0] & none & +8.3 [$-$27.5, +44.2] & none \\
\addlinespace
\multicolumn{6}{@{}l@{}}{\textit{Tier 3 + geometry + DFT magnetic outputs}} \\
\quad Gradient boosting & 205.6 & $-$2.8 [$-$39.1, +33.6] & none & +11.5 [$-$14.4, +37.4] & none \\
\quad XGBoost & 201.9 & +7.5 [$-$33.8, +48.8] & none & +17.2 [$-$9.5, +43.9] & borderline gain \\
\addlinespace
\multicolumn{6}{@{}l@{}}{\textit{Tier 3 + CHGNet crystal embedding (64)}} \\
\quad Gradient boosting & 213.7 & $-$10.8 [$-$22.5, +0.9] & none & +3.5 [$-$19.2, +26.2] & none \\
\quad XGBoost & 199.9 & +9.5 [$-$28.7, +47.7] & none & +19.2 [+0.4, +38.0] & stable gain \\
\bottomrule
\end{tabularx}
\par\vspace{4pt}\raggedright\noindent Matched Curie population (4,602 records), same 10 splits as Extended Data Table 1; MAE of Tier 1 / Tier 3: gradient boosting 202.9 / 217.1 K, XGBoost 209.3 / 219.1 K. Geometry: magnetic-ion site fraction and density, nearest and mean M-M distance, M-M connectivity, M-X distance and coordination, mean M-X-M angle and the fractions of angles in [150, 180] and [80, 100] degrees, from the Materials Project structure of the same entry. DFT magnetic outputs (moment per atom, magnetic-site fraction, magnetization per volume, magnetic flag) are not label leakage for experimental Curie temperatures. CHGNet = frozen 64-d crystal feature of the pretrained universal potential. Improvement in K, positive = better; Nadeau-Bengio 95\% intervals; verdict 'none' = no stable change; XGBoost single realization.
\end{table}
